\chapter{From Planner + Skills to VLA}

\section*{학습 목표}

이 장은 planner와 skill library를 결합한 구조가 왜 VLA로 이어졌는지 설명한다. Chapter 11에서는 LLM/VLM이 task를 분해하고, 기존 skill과 controller를 호출하는 구조를 다루었다. 그 구조는 해석 가능하고 안전장치를 붙이기 쉽지만, skill boundary가 너무 단단하면 새로운 행동을 만들기 어렵다. 반대로 end-to-end VLA policy는 perception, language, action을 하나의 mapping으로 학습할 수 있지만, control boundary와 safety 검증을 약화시킬 수 있다.

\begin{enumerate}
  \item planner+skills 구조의 engineering benefit과 bottleneck을 구분한다.
  \item skill boundary를 learned policy가 흡수할 때 생기는 장점과 위험을 설명한다.
  \item skill call, waypoint, action chunk, residual command를 VLA interface spectrum으로 정리한다.
  \item hierarchical VLA가 pure planning과 pure end-to-end 사이의 실용적 middle ground임을 이해한다.
\end{enumerate}

\section{Planner + skills 구조의 기본형}

Planner+skills 시스템은 다음 세 단계로 쓸 수 있다.
\begin{equation}
  \ell,\obs_t
  \xrightarrow{M_{\theta}^{plan}}
  P_t
  \xrightarrow{\mathrm{Select}}
  s_t
  \xrightarrow{\pi_{s_t}}
  \ctrlcmd_t.
  \label{eq:planner-skill-basic}
\end{equation}
여기서 $\ell$은 language instruction, $P_t$는 plan, $s_t$는 선택된 skill, $\pi_{s_t}$는 skill에 연결된 controller 또는 policy이다. 이 구조의 장점은 module boundary가 분명하다는 것이다. planner가 틀렸는지, skill precondition이 틀렸는지, controller가 실패했는지 분리해서 진단할 수 있다.

또한 이 구조는 data efficiency가 좋다. 로봇이 이미 grasp, place, navigate 같은 skill을 갖고 있으면 LLM/VLM planner는 새로운 task 조합을 비교적 적은 data로 만들 수 있다. 이때 학습 문제는 motor control 전체가 아니라 skill selection과 ordering으로 축소된다.
\begin{equation}
  s_t^{\star}
  =
  \arg\max_{s_i\in\mathcal{S}}
  Q_{\theta}(s_i \mid \ell,\bel_t,h_t),
  \label{eq:skill-selection-q}
\end{equation}
여기서 $h_t$는 history이다. 이 식에서 policy의 action space는 $\mathcal{S}$로 제한된다. 이는 장점이자 한계다.

\section{Skill boundary의 병목}

Skill library가 강하면 system은 안정적이다. 하지만 skill boundary가 task의 자연스러운 단위를 잘못 자르면 generalization이 막힌다. 예를 들어 ``잡기'' skill이 fixed top-down grasp만 지원한다면, 컵 손잡이, 얇은 천, 투명 용기, partially occluded object에 대해 planner가 아무리 좋은 sequence를 만들어도 실행이 어렵다.

Skill boundary의 병목은 세 가지로 나타난다.
\begin{align}
  \mathcal{B}_{coverage}
  &= \mathbb{I}[\exists s_i\in\mathcal{S}: s_i \text{ can express the required behavior}], \\
  \mathcal{B}_{parameter}
  &= \mathbb{I}[\exists \alpha_i: \pi_{s_i}(\alpha_i,\bel_t) \text{ succeeds}], \\
  \mathcal{B}_{feedback}
  &= \mathbb{I}[\text{skill exposes enough intermediate feedback}].
  \label{eq:skill-bottlenecks}
\end{align}
첫 번째는 필요한 행동이 skill set에 있는가의 문제다. 두 번째는 skill은 있지만 parameterization이 충분한가의 문제다. 세 번째는 skill 내부 상태가 planner나 monitor에 충분히 노출되는가의 문제다.

\begin{table}[h]
  \centering
  \caption{Skill boundary가 만드는 trade-off}
  \begin{tabularx}{\linewidth}{p{0.22\linewidth}X X}
    \toprule
    Boundary choice & 장점 & 병목 \\
    \midrule
    Coarse skill & 안전 검증과 reuse가 쉽다 & 새로운 contact mode와 세밀한 조작이 어렵다 \\
    Fine skill & 조합성이 좋아진다 & long-horizon planning과 error accumulation이 커진다 \\
    Parameterized skill & 기존 controller를 유지하면서 task variation을 처리한다 & parameter space가 커지면 planner가 사실상 control 문제를 푼다 \\
    Learned skill & data에서 motion pattern을 흡수한다 & 실패 원인과 safety certificate가 불명확해질 수 있다 \\
    \bottomrule
  \end{tabularx}
  \label{tab:skill-boundary-tradeoff}
\end{table}

Real VLA 관점에서 중요한 질문은 ``skill을 쓸 것인가 말 것인가''가 아니다. 어떤 boundary를 고정하고, 어떤 boundary를 학습할 것인가이다.

\section{VLA가 흡수하는 responsibility}

VLA는 language와 observation에서 robot action을 바로 예측하는 방향으로 skill boundary의 일부를 흡수한다.
\begin{equation}
  \act_t
  =
  \pi_{\theta}(\ell,\obs_{\leq t},\conf_t).
  \label{eq:vla-policy-direct}
\end{equation}
이 식은 planner+skills 구조와 다르게 skill name을 명시적으로 거치지 않는다. 대신 policy 내부 representation이 object, relation, subgoal, motion primitive를 latent하게 표현할 수 있다.

하지만 모든 responsibility가 하나의 neural network 안으로 사라지는 것은 아니다. 실제 system에서는 VLA output을 다음 중 하나의 interface로 제한한다.
\begin{equation}
  \act_t
  \in
  \{\mathrm{skill\ call},\mathrm{waypoint},\mathrm{delta\ pose},
  \mathrm{action\ chunk},\mathrm{joint\ command},\mathrm{residual}\}.
  \label{eq:vla-interface-set}
\end{equation}
각 interface는 control boundary를 다르게 만든다. skill call은 planner+skills에 가깝고, joint command는 end-to-end control에 가깝다. waypoint와 action chunk는 중간 지점이다.

\section{Interface spectrum}

Planner+skills에서 VLA로 넘어가는 과정을 하나의 spectrum으로 보자.
\begin{equation}
  \mathcal{I}
  =
  \{\mathcal{I}_{symbolic},
  \mathcal{I}_{skill},
  \mathcal{I}_{waypoint},
  \mathcal{I}_{chunk},
  \mathcal{I}_{lowlevel}\}.
  \label{eq:interface-spectrum}
\end{equation}

\begin{table}[h]
  \centering
  \caption{Planner+skills와 VLA 사이의 action interface}
  \begin{tabularx}{\linewidth}{p{0.2\linewidth}X X}
    \toprule
    Interface & Output & Engineering implication \\
    \midrule
    Symbolic plan & subgoal sequence, object relation & interpretability는 높지만 executability gap이 크다 \\
    Skill call & skill id, arguments & 안전하고 재사용 가능하나 skill coverage에 묶인다 \\
    Waypoint & end-effector pose, target pose & motion planner와 결합하기 쉽고 controller boundary가 남는다 \\
    Action chunk & short horizon action sequence & visuomotor smoothness를 학습하지만 latency와 stability를 관리해야 한다 \\
    Low-level command & joint velocity, torque, gripper command & 표현력은 크지만 safety와 transfer가 어렵다 \\
    \bottomrule
  \end{tabularx}
  \label{tab:interface-spectrum}
\end{table}

이 표는 VLA 논문을 읽을 때 가장 먼저 확인해야 하는 항목이다. 같은 VLA라는 이름을 쓰더라도 action interface가 다르면 학습 난이도, 필요한 data, safety checker, controller integration이 모두 달라진다.

\section{Latent skill과 hierarchical policy}

Planner+skills 구조에서 skill은 사람이 정의한 discrete object였다. VLA로 넘어가면 skill이 latent variable로 바뀔 수 있다.
\begin{align}
  z_t &\sim p_{\theta}(z_t \mid \ell,\obs_t,h_t), \\
  \act_{t:t+H} &\sim p_{\phi}(\act_{t:t+H} \mid z_t,\obs_t,\conf_t).
  \label{eq:latent-skill}
\end{align}
여기서 $z_t$는 explicit skill id가 아니라 learned latent skill이다. 이 구조는 skill library의 coverage 병목을 줄일 수 있다. data가 충분하면 latent skill은 사람이 이름 붙이지 않은 motion pattern도 표현할 수 있다.

하지만 latent skill은 검증이 어렵다. explicit skill은 precondition과 termination condition을 붙일 수 있지만, latent variable은 의미가 불명확할 수 있다. 따라서 hierarchical VLA는 보통 두 가지 장치를 함께 둔다.
\begin{equation}
  \begin{aligned}
  z_t &= H_{\theta}(\ell,\obs_t,h_t), \\
  \act_t &= L_{\phi}(z_t,\obs_t,\conf_t), \\
  \mathrm{allow}_t &= S_{\eta}(z_t,\act_t,\bel_t,\conf_t).
  \end{aligned}
  \label{eq:hierarchical-vla-safety}
\end{equation}
$H_{\theta}$는 high-level reasoning, $L_{\phi}$는 low-level action generation, $S_{\eta}$는 safety gate이다. 이 구조는 Chapter 18의 hierarchical VLA와 직접 연결된다.

\section{Residual VLA}

Planner+skills와 end-to-end VLA 사이의 실용적인 중간 형태는 residual command이다. 기존 controller가 nominal command를 만들고, learned policy가 보정항을 더한다.
\begin{equation}
  \ctrlcmd_t
  =
  \ctrlcmd_t^{base}
  +
  \Delta \ctrlcmd_t^{vla},
  \qquad
  \Delta \ctrlcmd_t^{vla}
  =
  f_{\theta}(\ell,\obs_t,\bel_t).
  \label{eq:residual-vla}
\end{equation}
이 구조는 classical control의 안정성을 유지하면서 perception-language 조건을 반영할 수 있다. 예를 들어 nominal trajectory는 motion planner가 만들고, VLA는 grasp approach angle, contact timing, local correction을 보정한다.

Residual 구조의 위험은 learned residual이 base controller의 안정성 가정을 깨는 것이다. 따라서 residual은 norm bound, frequency bound, safety projection을 가져야 한다.
\begin{equation}
  \Delta \ctrlcmd_t^{safe}
  =
  \Pi_{\mathcal{U}_{safe}}
  \left(\Delta \ctrlcmd_t^{vla}\right),
  \qquad
  \|\Delta \ctrlcmd_t^{safe}\| \leq \epsilon.
  \label{eq:residual-projection}
\end{equation}
이런 형태는 Real VLA에서 중요하다. learned model을 쓰되, controller와 safety boundary를 완전히 버리지 않는다.

\section{Closed-loop feedback의 위치 변화}

Planner+skills 구조에서는 feedback이 주로 planner 밖에서 처리된다. skill이 실패하면 monitor가 report를 만들고, planner가 다음 skill을 다시 고른다.
\begin{equation}
  h_{t+1}
  =
  \mathrm{Update}(h_t,r_t,e_t).
  \label{eq:external-feedback-update}
\end{equation}
VLA에서는 feedback의 일부가 policy 내부 observation history로 들어간다.
\begin{equation}
  \act_t
  =
  \pi_{\theta}(\ell,\obs_{t-L:t},\act_{t-L:t-1}).
  \label{eq:vla-feedback-history}
\end{equation}
이 변화는 중요하다. policy가 실패 징후를 직접 보고 local correction을 할 수 있기 때문이다. 그러나 history를 넣는 것만으로 recovery가 보장되지는 않는다. policy가 failure mode를 학습하려면 dataset에 실패, correction, recovery가 포함되어야 한다.

따라서 VLA data engine은 success trajectory만 모아서는 부족하다.
\begin{equation}
  \mathcal{D}
  =
  \mathcal{D}_{success}
  \cup
  \mathcal{D}_{failure}
  \cup
  \mathcal{D}_{correction}
  \cup
  \mathcal{D}_{intervention}.
  \label{eq:vla-data-mixture}
\end{equation}
이 식은 Chapter 15의 data engine과 연결된다. VLA가 planner+skills의 robustness를 넘어서려면 실패와 개입을 학습해야 한다.

\section{Middle ground로서의 hierarchical VLA}

Pure planning과 pure end-to-end policy는 양극단이다. 실제 로봇 시스템에서는 다음 objective가 동시에 필요하다.
\begin{equation}
  \max_{\theta,\phi}
  \quad
  J_{task}
  -
  \lambda_1 J_{unsafe}
  -
  \lambda_2 J_{latency}
  -
  \lambda_3 J_{uninterpretable}.
  \label{eq:vla-system-objective}
\end{equation}
이 objective에서 task success만 최적화하면 unsafe behavior가 생길 수 있고, interpretability만 중시하면 skill boundary가 너무 rigid해진다. hierarchical VLA는 이 trade-off를 조절한다.

\begin{tcolorbox}[title={Design principle: choose the learning boundary}]
Real VLA 설계에서 핵심 결정은 model size보다 learning boundary이다. 어떤 판단을 LLM/VLM planner에 맡길지, 어떤 motion을 learned action head에 맡길지, 어떤 constraint를 controller와 safety monitor에 고정할지 먼저 정해야 한다.
\end{tcolorbox}

\section{VLA로 넘어가는 조건}

Planner+skills에서 VLA로 넘어갈 필요가 생기는 조건은 비교적 분명하다.
\begin{enumerate}
  \item task variation이 skill library의 coverage를 넘어선다.
  \item perception과 action 사이의 fine-grained coupling이 중요하다.
  \item contact나 deformation처럼 symbolic subgoal로 표현하기 어려운 dynamics가 많다.
  \item long-horizon task에서 skill 사이의 error accumulation이 커진다.
  \item data가 충분해서 end-to-end 또는 hierarchical policy를 학습할 수 있다.
\end{enumerate}

반대로 모든 문제를 VLA로 바꿀 필요는 없다. stable pick-place, fixture가 잘 정의된 industrial task, safety-critical constrained motion은 explicit controller와 verified planner가 여전히 유리할 수 있다. Real VLA는 기존 robotics를 대체하는 이름이 아니라, 기존 robotics boundary를 학습 가능한 representation과 결합하는 설계 방식이다 \cite{saycan2022,codeaspolicies2022,rt12022}.

\chapterwork
{Planner+skills 구조와 RT-1 이후의 action-producing transformer 구조를 비교해서 읽어라 \cite{saycan2022,rt12022}.}
{skill library가 너무 강하면 generalization이 막히고, end-to-end policy가 너무 강하면 safety/control boundary가 약해지는 이유는 무엇인가?}
{기존 skill library가 있는 robot에 VLA policy를 추가한다고 가정하고, 어떤 boundary를 학습하고 어떤 boundary를 유지할지 표로 정리하라.}

\section{요약}

Planner+skills 구조는 VLA의 전 단계이면서 동시에 현대 VLA 시스템의 일부로 남아 있다. 이 구조는 해석 가능성, safety gate, controller reuse라는 장점이 있지만, skill coverage, parameterization, feedback exposure에서 병목을 만든다. VLA는 이 병목의 일부를 learned policy로 흡수한다. 그러나 모든 것을 end-to-end로 밀어 넣으면 system diagnosis와 safety assurance가 어려워진다.

따라서 planner+skills에서 VLA로의 전환은 단순한 model replacement가 아니다. action interface와 learning boundary를 재설계하는 문제다. 다음 장에서는 이 전환이 실제로 Robotics Transformer 계열에서 어떻게 나타났는지, robot action이 sequence prediction과 token prediction 문제로 재정의된 과정을 다룬다.
